\chapter{Introduction}
% Background Information: Introduce the field of motion tracking, its significance in various applications (e.g., sports, healthcare, robotics), and why AI is critical in advancing the field.
% Problem Statement: Clearly state the problem your research will address (e.g., limitations of current motion tracking models, inefficiencies in model generation, etc.).
% Research Objectives: Define your research objectives for this phase, such as identifying the challenges, understanding existing solutions, and defining the scope of your future framework.
% Motivation/Need for the Research: Explain why this research is needed, highlighting current gaps in the field and how your framework will help solve these issues.
% Research Questions: Pose key questions your research will answer (e.g., What challenges do existing models face? What makes AI a suitable approach for motion tracking?).
% Structure of the Report: Briefly describe the layout of the paper.

\section{Background}
% Motion tracking is a rapidly growing field that involves capturing and analyzing the movement of objects or individuals. It is a cornerstone technology in various domains, including healthcare, where it is used for physical rehabilitation and fall detection, sports for performance tracking and injury prevention, gaming for immersive VR/AR experiences, and robotics for navigation and control.

% The advent of advanced sensors and artificial intelligence (AI) has significantly enhanced the scope and precision of motion tracking systems. AI models, in particular, are critical in processing motion data, enabling applications such as gesture recognition, activity classification, and movement prediction. However, building these AI models involves complex processes like data collection, preprocessing, feature extraction, and model training. This complexity creates a barrier for developers, especially those without specialized expertise in AI or motion tracking.

% To address these challenges, there is a growing interest in creating frameworks that simplify AI model generation. Such frameworks can democratize motion tracking by making it accessible to a broader audience, including researchers, developers, and hobbyists.

Motion tracking is a dynamic and rapidly evolving field that involves the capture, processing, and analysis of motion, whether it is the movement of objects, individuals, or environments. According to \cite{article:human_motion_tracking_survey}, the technology plays a pivotal role in a variety of fields. In healthcare, motion tracking supports physical rehabilitation, aids in fall detection, and aids in patient monitoring. In sports, it enhances performance tracking, aids in injury prevention, and provides insights to athletes and coaches. In gaming and virtual reality (VR) or augmented reality (AR), it promotes immersive and interactive user experiences. Furthermore, in robotics, motion tracking supports navigation systems and enables precise control mechanisms.

Advances in sensor technology and artificial intelligence (AI) have been transformative, dramatically increasing the accuracy, efficiency, and flexibility of motion tracking systems. In particular, AI models are indispensable for interpreting motion data, leading to applications such as gesture recognition, activity classification, and predictive motion modeling. However, the process of developing these AI models is inherently complex. It involves several complex steps, including collecting motion data, preprocessing and filtering that data to remove noise, extracting meaningful features, and finally training machine learning or deep learning models. This high level of complexity poses a significant challenge, especially for developers and researchers who may lack expertise in AI or motion tracking technology.

To overcome these challenges and expand the accessibility of motion tracking technologies, there is an increasing focus on creating frameworks designed to simplify the creation of AI models. These frameworks aim to democratize the field, making motion tracking and AI integration more accessible to a wider audience, including not only seasoned experts and researchers, but also developers and enthusiasts with limited expertise. By providing streamlined tools and processes, such frameworks promise to lower barriers to entry and foster innovation.

\section{Problem Statement}
Despite advancements in motion tracking technologies, generating AI models for motion analysis remains a challenging task. Current approaches often require extensive technical knowledge, domain-specific expertise, and significant time and resources. Moreover, existing solutions are frequently tailored to specific applications, lacking the flexibility and scalability needed to adapt to diverse use cases.

% For instance, many available frameworks either demand a steep learning curve or provide limited functionality for customizing models to specific motion tracking requirements. This restricts innovation and hinders the widespread adoption of AI-powered motion tracking solutions.

% Thus, there is a need for a versatile, user-friendly framework that simplifies the generation of AI models for motion tracking, enabling both experts and non-experts to leverage its potential effectively.

For example, popular frameworks and tools often have steep learning curves, requiring users to invest significant effort to understand and utilize their capabilities. Additionally, these tools may offer limited functionality, limiting the ability to customize models for unique or emerging motion tracking requirements. Such limitations can stifle innovation, as they discourage experimentation and limit the adoption of AI-enabled motion tracking solutions in areas where they could have transformative impact.

There is therefore an urgent need for a flexible, intuitive, and user-friendly framework that simplifies the creation of AI models suitable for motion tracking applications. This framework must be designed to accommodate both experts and non-experts, empowering a diverse group of users to effectively exploit the potential of motion tracking technologies.

\section{Research Objectives}
% The primary objective of this research is to explore the feasibility and necessity of a framework for generating AI models tailored to motion tracking applications. Specifically, this research aims to:

% Investigate the challenges associated with creating AI models for motion tracking.
% Analyze existing frameworks and tools for AI model generation, such as Google’s Teachable Machine, to identify gaps and limitations.
% Propose the foundational design for a flexible and user-friendly framework that supports a variety of motion tracking applications.
% This research will provide the basis for the development of a comprehensive solution that simplifies AI model generation while maintaining performance and adaptability.

The overarching goal of this study is to evaluate the feasibility and importance of developing a specific framework for generating AI models in the context of motion tracking. To achieve this, the study will address the following specific objectives:

\begin{itemize}
\item To thoroughly investigate the inherent challenges associated with designing and deploying AI models for motion tracking applications.

\item To conduct a comprehensive analysis of existing frameworks and tools, such as Google's Teachable Machine, to identify current limitations and unmet needs.

\item To propose an initial conceptual design for a flexible, extensible, and user-friendly framework capable of supporting a wide range of motion tracking applications.

\item To explore the potential benefits of such a framework in enabling more efficient and accessible model development processes.

\end{itemize}

This research aims to lay the foundation for a solution that not only simplifies the technical aspects of AI model creation, but also ensures that the resulting models maintain high standards of performance and adaptability.

\section{Scope and Limitations}
% This research focuses on the intersection of AI and motion tracking, particularly on frameworks that facilitate AI model generation. The scope includes:

% Reviewing current methodologies and technologies for motion tracking.
% Evaluating the potential of sensor-based data collection and preprocessing pipelines.
% Identifying opportunities for designing a general-purpose AI model generator.
% However, the research is limited by certain factors. It will primarily address high-level design considerations rather than implementation details. Additionally, it will not delve deeply into domain-specific applications, such as medical diagnostics or sports analytics. Hardware and computational constraints may also restrict the scope of experimental validation.

This research focuses on the intersection of artificial intelligence and motion tracking, with a particular emphasis on frameworks that facilitate the creation of AI models. The scope includes the following key areas:

\begin{itemize}
\item A detailed review of current methodologies, technologies, and tools used in motion tracking.

\item An evaluation of sensor-based data collection techniques and pre-processing procedures to ensure data quality and reliability.

\item Identification of opportunities to design a general-purpose AI model generator that addresses both common and unique challenges in motion tracking.

\item Exploration of the technical design elements and user interfaces that would enable the widespread adoption of such a framework.

\end{itemize}

However, this research is subject to certain limitations. The focus will primarily be on high-level design and conceptual considerations rather than detailed implementation or coding. In addition, the study will not provide a comprehensive analysis of domain-specific applications, such as medical diagnostics or advanced sports analytics. Real-world testing may also be limited by hardware and computational constraints, which may impact the validation of proposed solutions.